\chapter{How does the world know which way to go?}\label{ch:research}

\section{No central calculator}
When light crosses from air into water, its route bends. The light does not pause at the boundary while an invisible planner surveys every possible trajectory. When warmth spreads from a hot object into a cooler one, no administrator compares all the molecules. The behaviour follows local physical relations, and those relations can have consequences that are described by compact global statements.

This is the feature we are looking for, not a claim that societies should behave like photons. No person could calculate the complete state of the planet each time a tap is opened or a delivery truck moves. If a proposed physical account required such an all-knowing observer, it would fail at the first ordinary decision. We need to understand how a local change can be assessed without pretending that the entire world is locally known.

Physics provides two kinds of discipline. Conservation tells us that a carrier has not appeared from nowhere inside a closed account. A relation of motion or change tells us which transitions occur under specified conditions. Neither is the same as a moral rule or a social allocation policy. Their value for our purpose is structural: they show how to make claims precise enough that hidden creation, omitted boundaries and impossible action can be detected.

\section{The bent path of light}
Imagine a point above a flat water surface and another below it. In a single uniform medium, the straight segment is the shortest geometrical distance. Across the boundary, light can instead follow two straight segments joined at an angle. The route is longer in space than the direct line, yet the light travels at different speeds in the two media. Distance alone is not the relevant quantity.

In elementary geometrical optics the travel time of a candidate path $\gamma$ is
\begin{equation}
 T[\gamma]=\frac{1}{c}\int_\gamma n(\mathbf r)\,ds ,
 \label{eq:fermat-time}
\end{equation}
where $c$ is the speed of light in vacuum, $n(\mathbf r)$ is the refractive index and $ds$ is a small piece of distance. The integral adds the time cost of all those pieces. Fermat's principle identifies the observed ray, in its domain of application, by stationary optical time: a small permitted displacement of the path produces no first-order change in $T$. The familiar phrase \emph{least time} is useful for the elementary two-medium case, but stationarity is the more general statement.\cite{intro-fermat}

We can see the rule in a drawing with no mystery. Let the light enter the surface at a horizontal coordinate $u$. If its starting point is a horizontal distance $a$ from the chosen origin and height $h_1$ above the surface, while its destination lies at horizontal coordinate $b$ and depth $h_2$, then the candidate time is
\begin{equation}
 T(u)=\frac{n_1}{c}\sqrt{(u-a)^2+h_1^2}
      +\frac{n_2}{c}\sqrt{(b-u)^2+h_2^2}.
\end{equation}
Differentiate with respect to the crossing point $u$. At a stationary path,
\begin{equation}
 \frac{dT}{du}=0
 \quad\Longrightarrow\quad
 n_1\sin\theta_1=n_2\sin\theta_2 ,
\end{equation}
which is Snell's law. The angles are measured from the normal to the interface. If $n_2>n_1$, the ray bends toward that normal. A single rule accounts for the bending because the cost of another small distance is different on the two sides.

\begin{workedbox}{Why the straight line can lose}
Take $n_1=1$ and $n_2=1.5$. If $\sin\theta_1=0.6$, stationarity requires $\sin\theta_2=0.4$. The path changes direction at the interface. It is not the shortest geometrical connection, because distance in the slower medium contributes more time.
\end{workedbox}

\section{What the analogy gives us, and what it does not}
The optical integral teaches us to name the quantity being accumulated and the path over which it is accumulated. It also teaches caution. Light is not awarded a payment for finding a route; Fermat's principle is not a doctrine that people should choose the fastest path or that an economy should minimise travel time. It is a physical result for light under specified optical conditions.

Later in this book another path integral will appear. It has a different mathematical character. If a smooth scalar potential $V$ is declared, integrating its gradient along a path between two states gives a difference of endpoint values. Such a gradient integral is path-independent while the optical time in Eq.~\eqref{eq:fermat-time} generally depends on which route is proposed. The similarity is the discipline of summing local contributions; the distinction is essential to the mathematics.

Nor should the expression \emph{the world knows} be taken literally. Local laws can generate orderly patterns without foresight. A photon does not inspect an economic objective, and molecules do not vote on the future. For human systems we must declare the actions, observations and decision rules ourselves. Physics will constrain and inform them, not silently select an institution.

\section{From path to balance}
A water movement between a reservoir and a treatment site also has a route. It may require a pump, electricity, maintenance and a pipe. Describing only the final delivery could hide the resource used along the way; counting the same transformation twice could invent an extra cost. The first task is therefore to keep the physical account closed: what left one place, what entered another, and what crossed the boundary?

But conservation is only the beginning. Even when an amount is conserved, useful energy can become dispersed; a system may remain energetically possible while losing its ability to perform a function. The next chapter examines the first and second laws of thermodynamics, then asks how a human body maintains function in a world where every activity consumes and transforms resources.
